\section{A permutation normal form with an exact memory charge}
\label{sec:r73-normal-form}

Theorem~\ref{thm:r72-reversible} identifies when finite memory exists. We now
determine the form of every possible finite realization, retaining the cost of
its prescribed pure initial label. The distinction between a minimal linear
representation and a positive realization remains essential. In particular,
the permutation representation carried by the observable space does not by
itself provide a decoder with values in the probability simplex.

Let $b$ be a compactly reversible response admitting a finite realization.
Write $G$ for its finite observable group, $d=\dim V_b$, and $N$ for its residual
cardinality. We use the chronological convention $g(u)g(v)=g(uv)$: the right
action on $V_b$ is $f\cdot g(w)=L_wf$, where $L_wf(z)=f(wz)$. Thus $G$ in this
notation is the opposite of the group of operators with the usual convention
for composition; inversion identifies the two groups. No group cardinality,
faithfulness, or representation-theoretic statement depends on this convention.
Every element has a representative in $A^*$, including its inverse, since a
finite group generated by the command elements is already generated by their
positive words. The value $b(g)=b(w)$ is consequently well defined.

A \emph{probability factorization on a finite $G$-set} consists of a finite set
$X$, row permutation matrices $P_g$ with $P_gP_h=P_{gh}$, a probability row
$\lambda$ on $X$, and a row-stochastic matrix $D$ with columns indexed by $C$,
such that
\begin{equation}\label{eq:r73-probability-factorization}
 b(g)=\lambda P_gD\qquad(g\in G).
\end{equation}
All rows of $D$ are probability laws, including rows associated with states
that are not distinguished by the immediate terminal measurement. Unused
orbits may be deleted. Define
\[
 \kappa_G(b)=\min\{|X|:\text{\eqref{eq:r73-probability-factorization} holds}\}.
\]
The residual realization shows that this minimum is finite and at most $N$.
It permits a mixed $\lambda$ without charging for its preparation; the
stochastic memory $m(b)$ continues to charge every live label, beginning at
the specified point mass as in \eqref{eq:r72-response}.

\begin{theorem}[Reversible stochastic normal form]\label{thm:r73-normal-form}
Every $M$-label stochastic realization of $b$ gives a faithful finite $G$-set
$X$ and a factorization \eqref{eq:r73-probability-factorization} satisfying
\begin{equation}\label{eq:r73-initialization-charge}
 |X|+\mathbf 1_{\{\lambda\text{ is not a point mass}\}}\le M.
\end{equation}
Conversely, every such factorization has a realization with the displayed
number of labels. Its transitions on $X$ are permutations; if $\lambda$ is
mixed, its only additional label initializes the distribution on the first
command. Thus arbitrary continuing stochastic dynamics confers no advantage
over one initialization mixture followed by a permutation action.
\end{theorem}

\begin{proof}
Fix a realization $(e_1^{\mathsf T},\{Q_a\},D_0)$, and retain the reachable
row space $R$ and the surjection $\pi:R\to V_b$ from the proof of
Theorem~\ref{thm:r72-reversible}. Write
\[
 \mathcal L=\{w\in A^*:g(w)=1\}.
\]
These are identity words in the observable group; their matrices need not be
identity matrices on hidden states. Choose finitely many words
$\ell_1,\ldots,\ell_s\in\mathcal L$ whose positive matrix entries cover all
directed edges occurring in any $Q_\ell$, $\ell\in\mathcal L$. There are at
most $M^2$ distinct edges, so $s\le M^2$ is possible. If no nontrivial choice
is needed, the empty word suffices. Set
\[
 T=\frac{I+Q_{\ell_1}+\cdots+Q_{\ell_s}}{s+1}.
\]
Its closed communicating classes $C_1,\ldots,C_k$ are closed for every
$Q_\ell$ with $\ell\in\mathcal L$. The positive diagonal makes the
restriction of $T$ to each such class aperiodic. The elementary convergence
argument in Lemma~\ref{lem:r73-lazy-convergence} gives a stochastic idempotent
\[
 E=\lim_{n\to\infty}T^n,
 \qquad
 \Delta_ME=\operatorname{conv}\{\mu_1,\ldots,\mu_k\},
\]
where $\mu_i$ is the unique stationary probability row of $T$ on $C_i$.
Their supports are disjoint, so this image is a simplex with exactly $k$
vertices. Every identity-word matrix preserves each $C_i$; hence
\begin{equation}\label{eq:r73-sandwich}
 EQ_\ell E=E\qquad(\ell\in\mathcal L).
\end{equation}
This is a sandwiched equality. We do not assert $Q_\ell E=E$, $EQ_\ell=E$,
or common stationarity of the $\mu_i$ under all identity words.

Each $T^n$ is a convex combination of identity-word matrices. Consequently
$RE\subseteq R$ and $\pi(rE)=\pi(r)$ for $r\in R$. For a word $u$, put
$F_u=EQ_uE$, considered as a stochastic map of $\Delta_ME$ to itself.
Choose a word $\bar u$ representing $g(u)^{-1}$. Replacing the middle $E$
by the convex combinations $T^n$ and using
$u\ell\bar u\in\mathcal L$ gives
\[
 F_uF_{\bar u}=E=F_{\bar u}F_u.
\]
Thus $F_u$ is an affine automorphism of that simplex. Such an automorphism
permutes its vertices, because both it and its inverse preserve the simplex
and therefore preserve extremality. If $g(u)=g(v)$, the same inverse word
shows $F_uF_{\bar u}=F_vF_{\bar u}=E$, so $F_u=F_v$. Next,
$EQ_uT^nQ_vE$ is a convex combination of $F_{u\ell v}$, all equal to
$F_{uv}$. Passing to the limit proves $F_uF_v=F_{uv}$. We have obtained a
permutation action $P_g$ on the vertices $X=\{\mu_1,\ldots,\mu_k\}$.

This action is faithful. If $F_u$ is the identity on $\Delta_ME$, then
$rEQ_uE=rE$ for all rows $r$. For $r\in R$, application of $\pi$ shows
$L_u\pi(r)=\pi(r)$. Surjectivity gives $L_u=I$ on $V_b$, and so $g(u)=1$.
Write $e_1^{\mathsf T}E=\sum_i\lambda_i\mu_i$ and set
$D(i,\cdot)=\mu_iD_0$. Intertwining on $R$ gives, for every word $w$,
\[
 \lambda P_{g(w)}D=e_1^{\mathsf T}EQ_wED_0=b(w),
\]
including the empty word. If the original initial label lies in $C_i$, its
row under $E$ is exactly $\mu_i$. A mixed $\lambda$ therefore requires that
the initial label lie outside the union of the closed classes. The $k$
nonempty disjoint classes then use at most $M-1$ labels. This proves
\eqref{eq:r73-initialization-charge}; a pure $\lambda$ only requires $k\le M$.

For the converse, when $\lambda$ is a point mass its support point is the
initial label in $X$. If it is mixed, adjoin a label $\star$ with decoder
$D(\star,\cdot)=\lambda D$. On command $a$, its transition row into $X$ is
$\lambda P_{g(a)}$, and no transition returns to $\star$. The rows on $X$
are those of $P_{g(a)}$. Induction gives \eqref{eq:r73-probability-factorization}
for every nonempty word, and the decoder at $\star$ gives the empty word.
The first input is applied during initialization, so it is not lost or
encoded in an uncounted preparatory clock.
\end{proof}

The general structure of nonnegative idempotents and bounded groups of
nonnegative matrices is classical; see Brown~\cite{Brown64},
Schwarz~\cite{Schwarz64}, and Flor~\cite{Flor69}. The proof above supplies the
additional steps required for an arbitrary semigroup realization: the
identity-word projection is constructed before passing to the observable
quotient, the sandwiched action is proved to factor through that quotient,
and the prescribed pure start is charged. It does not assume that the
original stochastic matrices themselves form a group.

\begin{theorem}[Exact memory formula]\label{thm:r73-memory-formula}
With randomized terminal readout allowed and the initial label prescribed
pure, the minimum stochastic memory is
\begin{equation}\label{eq:r73-memory-formula}
 m(b)=\min_{(X,\lambda,D)}
 \left(|X|+\mathbf1_{\{\lambda\text{ is not a point mass}\}}\right)
 =\min\{N,1+\kappa_G(b)\}.
\end{equation}
The minimum ranges over precisely the probability factorizations
\eqref{eq:r73-probability-factorization}, not arbitrary linear factorizations.
\end{theorem}
\begin{proof}
The first equality is both directions of Theorem~\ref{thm:r73-normal-form}.
A point-mass factorization gives a deterministic-transition realization on
$X$, so $|X|\ge N$ by Theorem~\ref{thm:r72-reversible}. A mixed factorization
costs at least $1+\kappa_G(b)$. Conversely, the residual machine has $N$
labels, and a factorization attaining $\kappa_G(b)$ costs at most
$1+\kappa_G(b)$. These inequalities give the second equality, also when a
minimum factorization happens to have a pure start.
\end{proof}

\begin{corollary}[Sharp group bounds]\label{cor:r73-group-bound}
An $M$-label realization of a compactly reversible response satisfies
\[
 |G|\le M!,\qquad
 M<N\ \Longrightarrow\ |G|\le(M-1)!.
\]
Both bounds are attained by families below. In particular, a stochastic
memory saving over the residual count requires a faithful permutation core
on at most $M-1$ points.
\end{corollary}
\begin{proof}
The core action in Theorem~\ref{thm:r73-normal-form} embeds $G$ into $S_k$.
If $M<N$, its start cannot be pure, since that would produce a deterministic
machine with at most $M$ labels. Thus $k\le M-1$ in this case.
\end{proof}

Let $\nu_G(V_b)$ be the least cardinality of a finite $G$-set whose real
permutation module contains $V_b$, and let $\mu(G)$ be the least cardinality
of a faithful finite $G$-set. All references to $V_b$ here use the right-action
convention already fixed. A probability factorization gives $V_b$ as a
subquotient of $\mathbb R^X$: take the cyclic span of $\lambda$, then its
observable quotient. Averaging an inner product over finite $G$ splits every
invariant subspace orthogonally, so this subquotient embeds in
$\mathbb R^X$. We obtain
\begin{equation}\label{eq:r73-module-bounds}
 d\le\nu_G(V_b)\le\kappa_G(b),\qquad
 m(b)\ge\max\{d,\mu(G),\min(N,1+\nu_G(V_b))\}.
\end{equation}
These are lower bounds, not an identification of a linear module with a
bounded decoder. Proposition~\ref{prop:r73-decoder-obstruction} makes the
distinction strict.

\begin{corollary}[Rank rigidity]\label{cor:r73-rank-rigidity}
For a finite compactly reversible response,
\[
 m(b)=d\quad\Longleftrightarrow\quad N=d.
\]
In that case every realization with $d$ labels has permutation transition
matrices. If $N>d$, then $m(b)\ge d+1$.
\end{corollary}
\begin{proof}
Every $M$-label realization has $d\le M$. If $M=d$, the core in
Theorem~\ref{thm:r73-normal-form} has $k\ge d$ by its linear rank, hence
$k=d$ and its start is pure. Thus $N\le d$, whereas $N\ge d$ because the
residuals span $V_b$. Conversely $N=d$ gives a deterministic machine with
$d$ labels. For the assertion about the original matrices, a $d$-label
realization has $R=\mathbb R^d$ and the observable map $\pi$ is an
isomorphism. An inverse word for $g(a)$ therefore gives a stochastic inverse
of $Q_a$ on all of $\mathbb R^d$. A stochastic map with a stochastic inverse
is an affine automorphism of the full simplex and hence a permutation.
\end{proof}

\subsection{Two commands and a factorial deterministic gap}
Let $k\ge3$, let $P_g$ be the natural row permutation action of $S_k$, choose
a nonconstant $\lambda\in\Delta_k$, and choose a nonconstant column
$t\in[0,1]^k$. Two commands are the $k$-cycle $(1\,2\,\cdots\,k)$ and the
transposition $(1\,2)$; these generate $S_k$. Supply the binary response
\begin{equation}\label{eq:r73-symmetric-response}
 b(g)=\bigl(\lambda P_gt,\,1-\lambda P_gt\bigr).
\end{equation}
If the distinct coordinate values of $\lambda$ have multiplicities
$n_1,\ldots,n_j$, put $N_\lambda=k!/(n_1!\cdots n_j!)$.

\begin{proposition}[Exact symmetric-group family]\label{prop:r73-symmetric-family}
The response \eqref{eq:r73-symmetric-response} has observable group $S_k$,
dimension $k$, residual cardinality $N_\lambda$, and
\[
 \kappa_{S_k}(b)=k,\qquad
 m(b)=\min\{N_\lambda,k+1\}.
\]
When the coordinates of $\lambda$ are pairwise distinct, the least stochastic
memory is $k+1$ and the least deterministic-transition memory is $k!$.
For a point-mass $\lambda$, both are $k$. If either $\lambda$ or $t$ is
constant, the response instead has one residual and one-label memory.
\end{proposition}
\begin{proof}
For any nonconstant vector, differences between it and its transposes yield
a nonzero multiple of some $e_i-e_j$. Conjugating by all permutations shows
that the span of its orbit differences is
$\mathbf1^\perp$. The orbit of $\lambda$ therefore spans $\mathbb R^k$,
because every row has mass one. The orbit of $t$ also spans
$\mathbb R^k$: its differences span $\mathbf1^\perp$, and its coordinate
sum is positive, since a nonconstant nonnegative vector is nonzero.
Consequently the map
\[
 r\longmapsto\bigl(h\longmapsto(rP_ht,\,r\mathbf1-rP_ht)\bigr)
\]
is injective on $\mathbb R^k$. It identifies the reachable row space with
$V_b$, and identifies the observable action with the faithful natural
permutation action. Its residuals are precisely the distinct permutations
of $\lambda$, giving $d=k$ and $N=N_\lambda$. The displayed factorization
and the rank lower bound give $\kappa=k$. The memory formula proves the
remaining assertion; it also shows $N_\lambda\ge k$. The two constant
cases are immediate from \eqref{eq:r73-symmetric-response}.
\end{proof}

This is an exact prepared quantum experiment in every dimension $k$: take
$\rho=\operatorname{diag}(\lambda)$, permutation unitaries for the commands,
and terminal effect $E=\operatorname{diag}(t)$. The construction already
lies in the diagonal, classical subalgebra. It demonstrates a stochastic
versus deterministic memory separation, and does not assert a quantum
advantage for this family. It also proves sharpness of
Corollary~\ref{cor:r73-group-bound}: the point-mass family has $M=k$ and
$|G|=k!$, while pairwise distinct coordinates have $M=k+1<N=k!$ for $k\ge3$.

For a concrete rational example, take $k=3$,
$\lambda=(1/2,1/3,1/6)$, and $t=(0,1/4,1)^{\mathsf T}$.
The six first-outcome values are
\[
 \left\{\frac14,\frac38,\frac7{24},\frac{13}{24},
                 \frac{11}{24},\frac7{12}\right\}.
\]
It has $d=3$, $N=6$, $\kappa=3$, and $m=4$.

\subsection{The decoder constraint can remove the entire saving}
The representation-theoretic lower bound in
\eqref{eq:r73-module-bounds} need not attain the positive factorization
minimum, even for the same group, rank, and residual count as this example.

\begin{proposition}[A strict bounded-decoder obstruction]
\label{prop:r73-decoder-obstruction}
Let $u=(1,-1,0)$ and let $S_3$ act by its natural permutation matrices. The
binary response with first coordinate
\begin{equation}\label{eq:r73-endpoint-response}
 \beta(g)=\frac12+\frac14 uP_gu^{\mathsf T}
\end{equation}
has $d=3$, $N=6$, and $\nu_{S_3}(V_b)=3$, but
\[
 \kappa_{S_3}(b)=m(b)=6.
\]
It is realized by a prepared qubit with
\[
 U_0=\begin{pmatrix}1&0\\0&-1\end{pmatrix},\qquad
 U_1=\frac12\begin{pmatrix}1&\sqrt3\\\sqrt3&-1\end{pmatrix},
 \qquad \rho=E=|+\rangle\langle+|.
\]
Thus the least stochastic memory can differ for responses having the same
observable group, module, and residual cardinality.
\end{proposition}
\begin{proof}
The orbit of $u$ consists of its six distinct permutations and spans the
two-dimensional zero-sum subspace. Pairing with all permutations of $u$
is nondegenerate on that subspace. Including the constant outcome sum gives
dimension three and six distinct residuals. This representation is the
natural permutation module $\mathbb R^3$, and the rank bound gives
$\nu=3$.

Write $s=(1\,2)$ and $r=(1\,2\,3)$. Directly from
\eqref{eq:r73-endpoint-response},
\[
 \beta(1)=1,\quad\beta(s)=0,\quad
 \beta(r)=\beta(r^2)=\frac14,
 \quad\beta(\tau)=\frac34
 \quad\text{for the other two transpositions }\tau.
\]
Suppose a probability factorization exists on at most five points. Its
action is faithful, by the same residual argument as in
Theorem~\ref{thm:r73-normal-form}. Transitive $S_3$-sets have sizes
$1,2,3,6$, since their stabilizers have orders $6,3,2,1$; the subgroups of
order two are conjugate and the subgroup of order three is unique. Thus
our faithful action has one natural three-point orbit and, at most, one
two-point sign orbit together with fixed points. Two three-point orbits
already exceed the budget, and actions on only fixed and sign orbits
annihilate the alternating subgroup.

Let $a$ be any point of positive initial mass and write $q$ for the
first-coordinate decoder. The endpoint values force $q(a)=1$ and
$q(a\cdot s)=0$. No fixed point can carry mass. Within the natural orbit,
the point fixed by $s$ cannot carry mass, and both of the points exchanged
by $s$ cannot carry positive mass simultaneously. Thus the initial mass in
that orbit is concentrated at one moved point $a$; denote its total by
$A$. If a sign orbit is present its mass is concentrated at a single point;
write this total as $B$. We have $A+B=1$. Every even permutation has decoder
one on the mass in the sign orbit, and every odd permutation has decoder
zero there.

The images $a\cdot r$ and $a\cdot r^2$ are the other two points in the
natural orbit. One is $a\cdot s$, with decoder zero, and the other's decoder
is some $t\in[0,1]$. The two equal values $\beta(r)=\beta(r^2)=1/4$ give
$B=1/4$ and $At=0$. Hence $A=3/4$ and $t=0$. Of the two transpositions
different from $s$, one fixes $a$ and the other sends it to the third point.
Their first-outcome values would therefore be $A$ and zero, whereas both
must be $3/4$. This contradiction rules out all factorizations on at most
five points. The regular six-point residual machine gives $\kappa\le6$,
so $\kappa=6$ and Theorem~\ref{thm:r73-memory-formula} gives $m=6$.

For the physical realization, the adjoint actions of $U_0$ and $U_1$ on
the real Bloch plane are two reflections whose product is a rotation by
$2\pi/3$. The initial and terminal Bloch vector is the $x$-axis vector.
The first reflection sends it to its negative, the other two reflections
have inner product $1/2$ with it, and the two nonidentity rotations have
inner product $-1/2$. The Born probability is one half of one plus that
inner product. It has exactly the displayed six group values, so it is
\eqref{eq:r73-endpoint-response} after the specified generator assignment.
\end{proof}

The cyclic partition formula of the preceding revision is recovered as a
special case in Proposition~\ref{prop:r73-cyclic-module}. For the general
finite group, however, both the start charge and the bounded-decoder
constraint belong in the exact minimization problem. Neither can be
removed by a rank calculation.
